\documentclass[11pt,a4paper]{article}
\usepackage[T1]{fontenc}
\usepackage{isabelle,isabellesym}
\usepackage{amssymb}

\usepackage{pdfsetup}

\urlstyle{rm}
\isabellestyle{it}

\begin{document}

\title{The modular-reduction layer of ML-DSA (FIPS 204)}
\author{Amar Akshat}
\maketitle

\begin{abstract}
FIPS 204 fixes the prime modulus $q = 8380417$ for ML-DSA. Deployed implementations reduce values
modulo $q$ with a small layer of routines that the standard does not itself spell out: Montgomery
reduction, a Barrett-style \emph{reduce32}, a conditional addition and their composition. This entry
specifies what each routine must satisfy at the integer level and proves the arithmetic core of
Montgomery reduction: if $T$ is congruent to $A \cdot \mathit{QINV}$ modulo $2^{32}$ and lies in the
signed 32-bit range, then $(A - T q) / 2^{32}$ is congruent to $A \cdot 2^{-32}$ modulo $q$ and lies
strictly between $-q$ and $q$.

Two of the specifications depart deliberately from the bounds documented in the reference
implementation, because those bounds are off by one at an endpoint. Montgomery reduction is specified
on a half-open input domain, since the inclusive upper endpoint returns exactly $q$ and so violates
the documented strict bound; and \emph{reduce32} is specified on its true reachable output window,
which is asymmetric. Both departures are noted where they occur.
\end{abstract}

\tableofcontents

\parindent 0pt\parskip 0.5ex

\input{session}

\bibliographystyle{abbrv}
\bibliography{root}

\end{document}
